A continuación se describen los estadísticos empleados para comparar dos muestras independientes. Estos incluyen tanto estadísticos paramétricos como no paramétricos, y permiten evaluar diferencias entre distribuciones desde distintos enfoques, como la comparación de medias, de rangos o de estimadores robustos de localización.

\subsection{t-Student (t)}
\label{subsec:t-student}

El clásico estadístico basado en la distribución t de Student, descrito por ejemplo en \cite{wackerly2008mathematical}, que se emplea para comparar la media de dos muestras independientes, asumiendo varianzas poblacionales idénticas y finitas. Sean $X_1,\dots,X_m \sim F_X$ y $Y_1,\dots,Y_n \sim F_Y$ muestras aleatorias i.i.d. con medias $\mu_X$ y $\mu_Y$, y varianza común $\sigma^2$.

Se plantean las hipótesis $H_0: \mu_X - \mu_Y = 0 \quad \text{vs} \quad H_1: \mu_X - \mu_Y \neq 0$, y se considera el estadístico dado por:
$$t = \frac{\bar X-\bar Y}{S_p \sqrt{\frac{1}{m}+\frac{1}{n}}}$$
donde $S_p^2$ es la varianza combinada definida como:
$$S_p^2 = \frac{(m-1)S_X^2+(n-1)S_Y^2}{m+n-2}$$

Bajo los supuestos de normalidad e igualdad de varianzas, el estadístico $t$ sigue exactamente una distribución t de Student con $m+n-2$ grados de libertad, a partir de la cual se calcula el p-valor bilateral.

\subsection{t de Welch}
\label{subsec:welch}

El test t de Welch es la alternativa robusta para dos muestras independientes cuando se vulnera el supuesto de homocedasticidad ($\sigma_X^2 \neq \sigma_Y^2$). Esta prueba, propuesta por Welch en 1947 \cite{welch1947}. El estadístico está dado por:
$$t_W = \frac{\bar X - \bar Y}{\sqrt{\frac{S_X^2}{m} + \frac{S_Y^2}{n}}}$$

Bajo supuestos de independencia y normalidad, $t_W$ se aproxima a una distribución t de Student donde los grados de libertad $\nu$ se estiman mediante la ecuación de Welch-Satterthwaite \cite{wiki_welch}. En la implementación computacional, el p-valor se calcula utilizando esta aproximación con los grados de libertad ajustados.

\subsection{Wilcoxon (W)}
\label{subsec:wilcoxon}

El test de rangos de Wilcoxon \cite{randles1979introduction} evalúa la hipótesis nula de igualdad de distribuciones ($H_0: F_X = F_Y$). Para que la prueba sea estrictamente libre de distribución , se asume que $F_X$ y $F_Y$ son distribuciones \textbf{absolutamente continuas}, garantizando teóricamente la ausencia de empates. 

Sea $Z_1,\dots,Z_{m+n}$ la muestra combinada y $R(X_i)$ el rango de $X_i$ en dicha muestra. El estadístico se define como:
$$W = \sum_{i=1}^{m} R(X_i)$$

Este test es equivalente al estadístico $U$ de Mann-Whitney, el cual cuenta el número de excedencias de $Y$ sobre $X$:
$$U = \sum_{i=1}^{m} \sum_{j=1}^{n} \mathbf{1}(X_i < Y_j)$$
Relacionados mediante la identidad $U = W - \frac{m(m+1)}{2}$. Bajo $H_0$ se tiene:
$$\mathbb{E}[U] = \frac{mn}{2}, \quad \text{Var}(U) = \frac{mn(m+n+1)}{12}$$

Asintóticamente, si $m, n \to \infty$ de manera que la proporción $\frac{m}{m+n} \to \lambda \in (0,1)$, el estadístico estandarizado converge a una normal estándar:
$$Z = \frac{U - \mathbb{E}[U]}{\sqrt{\text{Var}(U)}} \xrightarrow{\;\;d\;\;} N(0,1)$$

La cota de Berry-Esseen asegura que el error de esta aproximación es de orden $O(N^{-1/2})$ \cite{serfling1980approximation} donde $N=M+n$, lo que justifica el uso de la aproximación normal para los tamaños muestrales de este estudio, permitiendo aislar el desempeño de la prueba sin interferencia del error de aproximación.

\subsection{Intervalo de confianza Hodges-Lehmann (HL)}
\label{subsec:Hodges-Lehmann-IC}

Este método asume explícitamente el \textbf{modelo de locación continuo}, donde $F_Y(x) = F_X(x - \Delta)$ para algún parámetro de desplazamiento $\Delta \in \R$. Bajo este modelo, el estimador puntual de Hodges-Lehmann \cite{randles1979introduction} se define como la mediana de las $m \times n$ diferencias de Walsh:
$$\hat{\Delta}_{HL} = \text{median}\{Y_j - X_i : i=1,\dots,m,\; j=1,\dots,n\}$$

El intervalo de confianza al $1-\alpha$ para el parámetro poblacional $\Delta$ se construye invirtiendo la prueba de Wilcoxon-Mann-Whitney. Ordenando las $mn$ diferencias como $D_{(1)} \leq \cdots \leq D_{(mn)}$, el intervalo está dado por:
$$\left[ D_{(k)}, \, D_{(mn-k+1)} \right]$$
donde el índice $k$ corresponde al cuantil $\alpha/2$ de la distribución nula exacta del estadístico $U$.

\subsection{Test condicional de permutación (P)}
\label{subsec:test-permutacion}

El test condicional de permutación \cite{randles1979introduction} evalúa $H_0: F_X = F_Y$ basándose en el principio de \textbf{intercambiabilidad}. Si la hipótesis nula es cierta, la asignación de las $m+n$ etiquetas a los datos observados es puramente aleatoria y equiprobable.

Se emplearon los siguientes estadísticos de contraste $T(X,Y)$: diferencia de medias; $\bar X-\bar Y$, diferencia de medianas; $\text{median}(X)-\text{median}(Y)$ y diferencia de medias recortadas; $\bar X_{\epsilon} - \bar Y_{\epsilon}$

El p-valor es la proporción de permutaciones espaciales cuyo $|T_{perm}|$ supera o iguala a $|T_{obs}|$. Cabe advertir que, si existe heterocedasticidad o diferencias en la forma de las colas, las muestras dejan de ser intercambiables y el test pierde su propiedad de nivel exacto, inflando el error Tipo I al rechazar $H_0$ no por diferencias en localización, sino por pérdida de intercambiabilidad general.

\subsection{Kolmogorov-Smirnov (KS)}
\label{subsec:kolmogorov-smirnov}

El test de Kolmogorov-Smirnov para dos muestras evalúa la hipótesis nula $H_0: F_X = F_Y$ donde $X_1,\dots,X_m \sim F_X$ y $Y_1,\dots,Y_n \sim F_Y$ dos muestras aleatorias independientes absolutamente continuas. Bajo estas condiciones, el estadístico de Kolmogorov es libre de distribución bajo la hipótesis nula, el cual se define mediante $$D_{m,n} = \sup_{x \in \mathbb{R}} |F_{X,m}(x) - F_{Y,n}(x)|$$
donde $F_{X,m}(x)$ y $F_{Y,n}(x)$ son las funciones de distribución empíricas de cada muestra.

Para tamaños muestrales pequeños, la distribución de $D_{m,n}$ bajo $H_0$ se obtiene mediante combinatoria. Sin embargo, para muestras grandes, el cálculo se vuelve computacionalmente ineficiente. En su lugar, se emplea la distribución asintótica del estadístico escalado. Según los resultados clásicos descritos por ejemplo en Shao \cite{shao2003mathematical}, para tamaños $m$ y $n$ grandes, bajo $H_0$ la probabilidad acumulada asintótica está dada por:
$$\lim_{m,n \to \infty} P\left(\sqrt{\frac{mn}{m+n}}D_{m,n} \leq t \right) = \sum_{j=-\infty}^{\infty} (-1)^{j} e^{-2j^2t^2}, \quad t > 0.$$